\SourceLeaf{030}
\textit{in cuius peripheria sint $2m$ puncta, in quibus $T=0$, totidemque, in quibus $U=0$, et quidem ita, ut singula posteriora inter bina priorum iaceant.}

Sit summa omnium coëfficientium $A,B$ etc. $K,L,M$ positive acceptorum $=S$, accipiaturque $R$ simul $>S\BareRoot{2}$ et $>1$\HistoricalNote{9}: tum dico in circulo radio $R$ descripto ea, quae in theoremate enunciata sunt, necessario locum habere. Scilicet designato brevitatis gratia eo puncto huius circumferentiae, quod $\frac1m45$ gradibus ab ipsius concursu cum laeva parte axis distat, sive pro quo $\varphi=\frac1m45^\circ$, per (1); similiter eo puncto, quod $\frac3m45^\circ$ ab hoc concursu distat, sive pro quo $\varphi=\frac3m45^\circ$, per (3); porro eo, ubi $\varphi=\frac5m45^\circ$, per (5) etc. usque ad $(8m-1)$, quod $\frac{8m-1}m45$ gradibus ab illo concursu distat, si semper versus eandem partem progrederis, (aut $\frac1m45^\circ$ a parte opposita), ita ut omnino $4m$ puncta in peripheria habeantur, aequalibus intervallis dissita: iacebit inter $(8m-1)$ et (1) unum punctum, pro quo $T=0$; nec non sita erunt similia puncta singula inter (3) et (5); inter (7) et (9); inter (11) et (13) etc., quorum itaque multitudo $2m$; eodemque modo singula puncta, pro quibus $U=0$, iacebunt inter (1) et (3); inter (5) et (7); inter (9) et (11), quorum multitudo igitur etiam $=2m$; denique praeter haec $4m$ puncta alia in tota peripheria non dabuntur, pro quibus vel $T$ vel $U$ sit $=0$.

\textit{Demonstr.} I. In puncto (1) erit $m\varphi=45^\circ$ adeoque
\[ T=R^{m-1}(R\BareRoot{\tfrac12}+A\sin(m-1)\varphi+\tfrac BR\sin(m-2)\varphi+\text{etc.}+\tfrac L{R^{m-2}}\sin\varphi)
\] summa vero $A\sin(m-1)\varphi+\frac BR\sin(m-2)\varphi$ etc. certo non poterit esse maior quam $S$, adeoque necessario erit minor quam $R\BareRoot{\frac12}$: unde sequitur, in hoc puncto valorem ipsius $T$ certo esse positivum. A potiori itaque $T$ valorem positivum habebit, quando $m\varphi$ inter $45^\circ$ et $135^\circ$ iacet, i. e. a puncto (1) usque ad (3) valor ipsius $T$ semper positivus erit. Ex eadem ratione $T$ a puncto (9) usque ad (11) positivum valorem ubique habebit, et generaliter a quovis puncto $(8k+1)$ usque ad $(8k+3)$, denotante $k$ integrum quemcunque. Simili modo $T$ ubique inter (5) et (7), inter (13) et (15) etc. et generaliter inter $(8k+5)$ et $(8k+7)$ valorem negativum habebit, adeoque in omnibus his intervallis nullibi poterit esse $=0$. Sed quoniam in (3) hic valor est positivus, in (5) negativus: necessario alicubi inter (3) et (5) erit $=0$; nec non alicubi inter (7) et (9); inter (11) et (13) etc.
\SourceLeaf{031}
usque ad intervallum inter $(8m-1)$ et (1) incl., ita ut omnino in $2m$ punctis habeatur $T=0$. Q. E. P.

II. Quod vero praeter haec $2m$ puncta, alia, hac proprietate praedita, non dantur, ita cognoscitur. Quum inter (1) et (3); inter (5) et (7) etc. nulla sint, aliter fieri non posset, ut plura talia puncta exstent, quam si in aliquo intervallo inter (3) et (5), vel inter (7) et (9) etc. ad minimum duo iacerent. Tum vero necessario in eodem intervallo $T$ alicubi esset \textit{maximum}, vel \textit{minimum}, adeoque $\frac{\mathrm dT}{\mathrm d\varphi}=0$. Sed $\frac{\mathrm dT}{\mathrm d\varphi}=mR^{m-2}(R\cos m\varphi+\frac{m-1}mA\cos(m-1)\varphi+\text{etc.})$ et $\cos m\varphi$ inter (3) et (5) semper est negativus et $>\BareRoot{\frac12}$. Unde facile perspicitur, in toto hoc intervallo $\frac{\mathrm dT}{\mathrm d\varphi}$ esse quantitatem negativam; eodemque modo inter (7) et (9) ubique positivam; inter (11) et (13) negativam etc., ita ut in nullo horum intervallorum esse possit 0, adeoque suppositio consistere nequeat. Quare etc. Q. E. S.

III. Prorsus simili modo demonstratur, $U$ habere valorem negativum ubique inter (3) et (5), inter (11) et (13) etc. et generaliter inter $(8k+3)$ et $(8k+5)$; positivum vero inter (7) et (9), inter (15) et (17) etc. et generaliter inter $(8k+7)$ et $(8k+9)$. Hinc statim sequitur, $U=0$ fieri debere alicubi inter (1) et (3), inter (5) et (7) etc., i. e. in $2m$ punctis. In nullo vero horum intervallorum fieri poterit $\frac{\mathrm dU}{\mathrm d\varphi}=0$ (quod facile simili modo ut supra probatur): quamobrem plura quam illa $2m$ puncta in circuli peripheria non dabuntur, in quibus fiat $U=0$. Q. E. T. et Q.

Ceterum ea theorematis pars, secundum quam plura quam $2m$ puncta non dantur, in quibus $T=0$, neque plura quam $2m$, in quibus $U=0$, etiam inde demonstrari potest, quod per aequationes $T=0$, $U=0$ exhibentur curvae $m^{\mathrm{ti}}$ ordinis, quales a circulo tamquam curva secundi ordinis in pluribus quam $2m$ punctis secari non posse, ex geometria sublimiori constat.

\Article{20} Si circulus alius radio maiori quam $R$ ex eodem centro describitur, eodemque modo dividitur: etiam in hoc inter puncta (3) et (5) iacebit punctum unum, in quo $T=0$, itemque inter (7) et (9) etc., perspicieturque facile, quo minus radius huius circuli a radio $R$ differat, eo propius huiusmodi puncta inter (3) et (5) in utriusque circumferentia sita esse debere. Idem etiam locum habebit, si circulus radio aliquantum minori quam $R$, attamen maiori quam $S\BareRoot{2}$ et 1, describitur. Ex his nullo negotio intelligitur, circuli radio $R$ descripti
